\documentclass[a4paper,11pt]{article}
\usepackage{lmodern}
\usepackage[T1]{fontenc}
\usepackage{amsmath,amssymb,amsthm}
\usepackage{longtable,booktabs,array,calc}
\usepackage{hyperref}
\usepackage[margin=25mm]{geometry}

\newtheorem{theorem}{Theorem}[section]
\newtheorem{proposition}[theorem]{Proposition}
\newtheorem{lemma}[theorem]{Lemma}
\newtheorem{corollary}[theorem]{Corollary}
\newtheorem{definition}[theorem]{Definition}
\newtheorem{remark}[theorem]{Remark}
\newtheorem{observation}[theorem]{Observation}

\providecommand{\tightlist}{\setlength{\itemsep}{0pt}\setlength{\parskip}{0pt}}
\newcounter{none}

\begin{document}

\section{The sine-Gordon Equation: Foundations of Topological Soliton
Theory}\label{the-sine-gordon-equation-foundations-of-topological-soliton-theory}

\subsection{-- Classification of Topological Solutions and State
Transformation Properties in Nonlinear Wave Equations
--}\label{classification-of-topological-solutions-and-state-transformation-properties-in-nonlinear-wave-equations}

\textbf{Author:} Noriaki Kihara (WF System Co., Ltd.)

\textbf{Date:} April 2026

\textbf{DOI:} 10.5281/zenodo.19650966

\begin{center}\rule{0.5\linewidth}{0.5pt}\end{center}

\subsection{0. Purpose of This Paper}\label{purpose-of-this-paper}

This paper provides an overview of the mathematical structure of the
\textbf{sine-Gordon equation}, one of the nonlinear wave equations. The
sine-Gordon equation is the most fundamental nonlinear equation that
admits topological solitons as solutions, and it is widely applied in
particle physics, condensed matter physics, and field theory.

In this paper, we organize the definition of the sine-Gordon equation,
the classification of solutions (kinks, anti-kinks, breathers), phase
shifts during collisions, and physical applications, tracing back to the
original papers. Furthermore, we prove that soliton solutions are
rigorously characterized by a finite number of parameters, and that all
structural properties of solitons are preserved in the limit of the
spread parameter \(\sigma \to 0\) (rectangular wave limit). In addition,
we define the transformation invariants of the sinusoidal-rectangular
wave transformation (§11), explicitly state the limits of applicability
of sine-Gordon theory (§9.8), and present an elastic state
transformation model on phase space that preserves the transformation
invariants as an alternative model that overcomes these limitations
(§12).

\begin{center}\rule{0.5\linewidth}{0.5pt}\end{center}

\subsection{1. Definition of the sine-Gordon
Equation}\label{definition-of-the-sine-gordon-equation}

\subsubsection{1.1 The Equation}\label{the-equation}

The 1+1-dimensional sine-Gordon equation is defined as follows {[}1,
2{]}:

\[\frac{\partial^2 \phi}{\partial t^2} - \frac{\partial^2 \phi}{\partial x^2} + \sin(\phi) = 0\]

Here \(\phi = \phi(x, t)\) is a scalar field (phase field), and \(x, t\)
are dimensionless independent variables. In this paper, we treat the
above equation as a dimensionless phase equation, and do not restore
physical unit systems.

\subsubsection{1.2 Origin of the Name}\label{origin-of-the-name}

The name ``sine-Gordon equation'' derives from its similarity to the
Klein-Gordon equation

\[\frac{\partial^2 \phi}{\partial t^2} - \frac{\partial^2 \phi}{\partial x^2} + m^2 \phi = 0\]

It is obtained by replacing the linear term \(m^2 \phi\) with the
nonlinear term \(\sin(\phi)\), and the name is a portmanteau combining
``sine'' and ``Gordon.'' This name was introduced by Rubinstein {[}3{]}.

\subsubsection{1.3 Lagrangian Density}\label{lagrangian-density}

The sine-Gordon equation is derived from the following Lagrangian
density {[}4{]}:

\[\mathcal{L} = \frac{1}{2}\left(\frac{\partial \phi}{\partial t}\right)^2 - \frac{1}{2}\left(\frac{\partial \phi}{\partial x}\right)^2 - (1 - \cos\phi)\]

The potential \(V(\phi) = 1 - \cos\phi\) attains its minimum value of 0
at \(\phi = 2n\pi\) (\(n \in \mathbb{Z}\)). The existence of these
infinitely many degenerate vacua is the origin of topological solitons.

\subsubsection{1.4 Historical Background}\label{historical-background}

The sine-Gordon equation was independently introduced in the following
contexts:

\begin{itemize}
\tightlist
\item
  \textbf{Crystal dislocation theory} (1939): Frenkel and Kontorova
  introduced it as a model describing the motion of dislocations in
  crystal lattices {[}1{]}.
\item
  \textbf{Differential geometry} (19th century): In the theory of
  surfaces with constant negative curvature (pseudospheres), Bour
  {[}5{]} and Enneper {[}6{]} had derived an isomorphic equation.
  Bianchi {[}7{]} studied the Bäcklund transformation of this equation.
\item
  \textbf{Particle physics} (1960s): Skyrme used a sine-Gordon type
  equation as a nonlinear field theory model for baryons {[}8{]}.
  Perring and Skyrme were the first to perform numerical calculations of
  kink-anti-kink collisions {[}9{]}.
\end{itemize}

\begin{center}\rule{0.5\linewidth}{0.5pt}\end{center}

\subsection{2. Classification of
Solutions}\label{classification-of-solutions}

\subsubsection{\texorpdfstring{2.1 Kink Solution (Winding Number
\(+1\))}{2.1 Kink Solution (Winding Number +1)}}\label{kink-solution-winding-number-1}

The kink solution is a localized wave in which \(\phi\) smoothly
transitions from \(0\) to \(2\pi\) {[}2, 9{]}:

\[\phi_K(x, t) = 4\arctan\left(\exp\left(\frac{x - vt}{\sqrt{1 - v^2}}\right)\right)\]

Here \(v\) (\(|v| < 1\)) is the velocity of the kink.

\textbf{Properties:}

\begin{itemize}
\tightlist
\item
  \(x \to -\infty\): \(\phi \to 0\); \(x \to +\infty\):
  \(\phi \to 2\pi\)
\item
  The phase increases by \(2\pi\) -- \textbf{winding number (topological
  charge)} \(N = +1\)
\item
  The denominator \(\sqrt{1 - v^2}\) is the Lorentz factor, covariant
  under special relativity
\item
  Rest energy \(E_0 = 8m\) (in natural units \(E_0 = 8\))
\end{itemize}

\subsubsection{\texorpdfstring{2.2 Anti-kink Solution (Winding Number
\(-1\))}{2.2 Anti-kink Solution (Winding Number -1)}}\label{anti-kink-solution-winding-number--1}

The anti-kink solution is a localized wave in which \(\phi\) transitions
from \(2\pi\) to \(0\) {[}2{]}:

\[\phi_{\bar{K}}(x, t) = 4\arctan\left(\exp\left(-\frac{x - vt}{\sqrt{1 - v^2}}\right)\right)\]

\textbf{Properties:}

\begin{itemize}
\tightlist
\item
  \(x \to -\infty\): \(\phi \to 2\pi\); \(x \to +\infty\):
  \(\phi \to 0\)
\item
  The phase decreases by \(2\pi\) -- \textbf{winding number} \(N = -1\)
\item
  Has the same mass and velocity characteristics as the kink
\end{itemize}

\subsubsection{2.3 Definition and Conservation of Winding
Number}\label{definition-and-conservation-of-winding-number}

\textbf{Definition 2.1.} The \textbf{winding number} (topological
charge) of a field configuration \(\phi(x, t)\) is defined as follows
{[}4, 10{]}:

\[N = \frac{1}{2\pi}\int_{-\infty}^{+\infty} \frac{\partial \phi}{\partial x}\,dx = \frac{\phi(+\infty) - \phi(-\infty)}{2\pi}\]

\textbf{Proposition 2.1 (Conservation of winding number).} Under the
time evolution of the sine-Gordon equation, the winding number \(N\) is
conserved.

\emph{Proof.} The boundary values \(\phi(\pm\infty)\) are fixed at
vacuum values (\(2n\pi\)), and cannot change by integer multiples of
\(2\pi\) through continuous time evolution. Therefore \(N\) is conserved
as a time-independent integer value. \(\square\)

This conservation law guarantees the ``particle-like nature'' of
solitons -- creation and annihilation are possible only in pairs (kink +
anti-kink).

\subsubsection{2.4 Breather Solution (Bound Oscillating
State)}\label{breather-solution-bound-oscillating-state}

A breather is an oscillating solution in which a kink and an anti-kink
are bound, with winding number \(N = 0\) {[}2, 11{]}:

\[\phi_B(x, t) = 4\arctan\left(\frac{\sqrt{1 - \omega^2}\;\cos(\omega t)}{\omega\;\cosh\left(\sqrt{1 - \omega^2}\;x\right)}\right)\]

Here \(\omega\) (\(0 < \omega < 1\)) is the internal oscillation
frequency of the breather.

\textbf{Properties:}

\begin{itemize}
\tightlist
\item
  A localized oscillating solution that decays exponentially in space
\item
  The smaller the internal frequency \(\omega\), the larger the spatial
  extent of the breather
\item
  The rest mass of the breather is \(M_B = 16\sqrt{1 - \omega^2}\)
\item
  In the limit \(\omega \to 0\), the breather mass approaches that of a
  kink-anti-kink pair (\(2 \times 8 = 16\))
\item
  In the limit \(\omega \to 1\), the breather mass approaches \(0\),
  transitioning to linear small oscillations (mesons)
\end{itemize}

\subsubsection{2.5 Summary of Solution
Classification}\label{summary-of-solution-classification}

{\def\LTcaptype{none} % do not increment counter
\begin{longtable}[]{@{}
  >{\centering\arraybackslash}p{(\linewidth - 8\tabcolsep) * \real{0.2381}}
  >{\centering\arraybackslash}p{(\linewidth - 8\tabcolsep) * \real{0.2857}}
  >{\centering\arraybackslash}p{(\linewidth - 8\tabcolsep) * \real{0.1905}}
  >{\centering\arraybackslash}p{(\linewidth - 8\tabcolsep) * \real{0.1429}}
  >{\centering\arraybackslash}p{(\linewidth - 8\tabcolsep) * \real{0.1429}}@{}}
\toprule\noalign{}
\begin{minipage}[b]{\linewidth}\centering
Solution Type
\end{minipage} & \begin{minipage}[b]{\linewidth}\centering
Winding Number \(N\)
\end{minipage} & \begin{minipage}[b]{\linewidth}\centering
Localization
\end{minipage} & \begin{minipage}[b]{\linewidth}\centering
Oscillation
\end{minipage} & \begin{minipage}[b]{\linewidth}\centering
Mass
\end{minipage} \\
\midrule\noalign{}
\endhead
\bottomrule\noalign{}
\endlastfoot
Kink & \(+1\) & Localized & None & \(8\) \\
Anti-kink & \(-1\) & Localized & None & \(8\) \\
Breather & \(0\) & Localized & Yes & \(16\sqrt{1-\omega^2}\) \\
Meson (linear wave) & \(0\) & Non-localized & Yes & \(\geq 2\)
(threshold) \\
\end{longtable}
}

\begin{center}\rule{0.5\linewidth}{0.5pt}\end{center}

\subsection{3. Soliton Collisions and Phase
Shifts}\label{soliton-collisions-and-phase-shifts}

\subsubsection{3.1 Kink-Kink Collision}\label{kink-kink-collision}

The collision between kinks of the same sign (\(N = +1\) and \(N = +1\))
is repulsive, and the two-body solution is given by {[}2, 12{]}:

\[\phi_{KK}(x, t) = 4\arctan\left(\frac{v\,\sinh\left(\frac{x}{\sqrt{1-v^2}}\right)}{\cosh\left(\frac{vt}{\sqrt{1-v^2}}\right)}\right)\]

After the collision, the two kinks separate again, completely preserving
their shapes. However, compared to the case without collision, the
position of each kink is shifted by a \textbf{phase shift} \(\Delta x\):

\[\Delta x = \frac{\sqrt{1 - v^2}}{v}\ln(v^2)\]

\subsubsection{3.2 Kink-Anti-kink
Collision}\label{kink-anti-kink-collision}

The collision between a kink and an anti-kink of opposite signs
(\(N = +1\) and \(N = -1\)) is attractive, and the two-body solution is
given by {[}2, 9, 12{]}:

\[\phi_{K\bar{K}}(x, t) = 4\arctan\left(\frac{v^{-1}\,\cosh\left(\frac{vt}{\sqrt{1-v^2}}\right)}{\sinh\left(\frac{x}{\sqrt{1-v^2}}\right)}\right)\]

\textbf{Properties:}

\begin{itemize}
\tightlist
\item
  With sufficient energy, the kink and anti-kink pass through each other
  and separate again (perfectly elastic collision: velocity and shape
  are identical before and after, only a phase shift occurs)
\item
  With insufficient energy, they form a breather (bound state)
  (oscillating permanently without decay)
\end{itemize}

\textbf{Remark:} In the sine-Gordon equation as a classical completely
integrable system, the annihilation of a kink-anti-kink pair followed by
conversion to radiation waves (mesons) does not occur. This is a
phenomenon that arises only through quantum effects or in systems where
integrability is broken (perturbation, discretization,
higher-dimensionalization).

\subsubsection{3.3 Physical Meaning of the Phase
Shift}\label{physical-meaning-of-the-phase-shift}

\textbf{Proposition 3.1 (Shape preservation and phase shift).} In
soliton collisions of the sine-Gordon equation:

\begin{enumerate}
\def\labelenumi{\arabic{enumi}.}
\tightlist
\item
  The \textbf{shape} of each soliton after the collision is identical to
  that before the collision.
\item
  The \textbf{position} of each soliton after the collision is displaced
  by a phase shift \(\Delta x\) compared to the case without collision.
\item
  The total winding number is conserved before and after the collision.
\end{enumerate}

\emph{Proof.} The sine-Gordon equation can be solved exactly by the
inverse scattering method {[}12{]}, and multi-soliton solutions
asymptotically separate into superpositions of individual solitons. The
asymptotic shape of each soliton is identical to that before the
collision, but with an added phase shift. Since the winding number is a
topological invariant, it is conserved. \(\square\)

This property is critically important for the present theory. The
\textbf{type} of soliton (= particle) (winding number = orientational
structure) is invariant under interactions, and only the
\textbf{position} (phase) changes.

\begin{center}\rule{0.5\linewidth}{0.5pt}\end{center}

\subsection{4. Complete Integrability and the Inverse Scattering
Method}\label{complete-integrability-and-the-inverse-scattering-method}

\subsubsection{4.1 Lax Pair}\label{lax-pair}

The sine-Gordon equation is a \textbf{completely integrable system} and
possesses a Lax pair {[}12, 13{]}.

Using light-cone coordinates \(u = \frac{1}{2}(x - t)\),
\(v = \frac{1}{2}(x + t)\), the equation becomes

\[\frac{\partial^2 \phi}{\partial u\,\partial v} = \sin\phi\]

This equation is expressed as the compatibility condition of the
following Lax pair:

\[\Psi_u = U\Psi, \quad \Psi_v = V\Psi\]

Here \(U\) and \(V\) are \(2 \times 2\) matrices depending on a spectral
parameter \(\lambda\) {[}12{]}.

\subsubsection{4.2 Infinitely Many Conserved
Quantities}\label{infinitely-many-conserved-quantities}

As a consequence of complete integrability, the sine-Gordon equation
possesses infinitely many independent conserved quantities {[}12, 14{]}.
The first three are as follows:

\begin{itemize}
\item
  \textbf{Energy:}
  \[E = \int_{-\infty}^{+\infty}\left[\frac{1}{2}\left(\frac{\partial\phi}{\partial t}\right)^2 + \frac{1}{2}\left(\frac{\partial\phi}{\partial x}\right)^2 + (1-\cos\phi)\right]dx\]
\item
  \textbf{Momentum:}
  \[P = -\int_{-\infty}^{+\infty}\frac{\partial\phi}{\partial t}\frac{\partial\phi}{\partial x}\,dx\]
\item
  \textbf{Winding number:}
  \[N = \frac{1}{2\pi}\int_{-\infty}^{+\infty}\frac{\partial\phi}{\partial x}\,dx\]
\end{itemize}

\subsubsection{4.3 Bäcklund
Transformation}\label{buxe4cklund-transformation}

The sine-Gordon equation possesses a Bäcklund transformation, which can
generate new solutions from known ones {[}7, 15{]}.

Given a known solution \(\phi_0\), the solution \(\phi_1\) of the
following system of equations is also a solution of the sine-Gordon
equation:

\[\frac{\partial}{\partial u}\left(\frac{\phi_1 + \phi_0}{2}\right) = a\,\sin\left(\frac{\phi_1 - \phi_0}{2}\right)\]

\[\frac{\partial}{\partial v}\left(\frac{\phi_1 - \phi_0}{2}\right) = \frac{1}{a}\sin\left(\frac{\phi_1 + \phi_0}{2}\right)\]

Here \(a\) is the parameter of the Bäcklund transformation. Starting
from the vacuum solution \(\phi_0 = 0\), one Bäcklund transformation
yields the kink solution, and two transformations yield the kink-kink
solution or breather solution {[}15{]}.

\begin{center}\rule{0.5\linewidth}{0.5pt}\end{center}

\subsection{5. Extension to Higher
Dimensions}\label{extension-to-higher-dimensions}

\subsubsection{5.1 2+1 and 3+1 Dimensions}\label{and-31-dimensions}

The direct extension of the sine-Gordon equation to higher dimensions is

\[\frac{\partial^2\phi}{\partial t^2} - \nabla^2\phi + \sin(\phi) = 0\]

However, in 2+1 dimensions and above, this equation is no longer
completely integrable {[}16{]}. Furthermore, by Derrick's theorem
{[}17{]}, stable localized solutions (solitons) consisting only of
scalar fields do not exist in 3+1 dimensions.

\subsubsection{5.2 Methods of
Stabilization}\label{methods-of-stabilization}

The following methods are known for obtaining stable solitons in higher
dimensions {[}10, 18{]}:

\begin{enumerate}
\def\labelenumi{\arabic{enumi}.}
\tightlist
\item
  \textbf{Skyrme model}: Adding higher-order derivative terms (Skyrme
  term) {[}8, 18{]}.
\item
  \textbf{Coupling with gauge fields}: Coupling a gauge field to the
  scalar field.
\item
  \textbf{Topological stabilization}: When the field carries a
  non-trivial topological charge, energy minimization prohibits the
  collapse of the soliton {[}10{]}.
\end{enumerate}

\subsubsection{5.3 Stabilization on Closed
Spaces}\label{stabilization-on-closed-spaces}

When dealing with solitons on higher-dimensional spaces, a mechanism to
circumvent Derrick's theorem is needed. On a closed curve \(S^1\) of
circumference \(L\), the premises of Derrick's theorem (asymptotic
conditions on infinite space) break down, so stable soliton solutions
can exist {[}10{]}.

\begin{center}\rule{0.5\linewidth}{0.5pt}\end{center}

\subsection{6. Physical Applications}\label{physical-applications}

\subsubsection{6.1 Condensed Matter
Physics}\label{condensed-matter-physics}

\begin{itemize}
\tightlist
\item
  \textbf{Josephson junctions}: The motion of magnetic flux quanta in
  tunnel junctions between superconductors is described by the
  sine-Gordon equation {[}19{]}. Kinks correspond to magnetic flux
  quanta (fluxons).
\item
  \textbf{Crystal dislocations}: Dislocations in crystal lattices in the
  Frenkel-Kontorova model correspond to kink solutions {[}1{]}.
\end{itemize}

\subsubsection{6.2 Particle Physics}\label{particle-physics}

\begin{itemize}
\tightlist
\item
  \textbf{Skyrme model}: Baryons are described as topological solitons
  (skyrmions) {[}8, 18{]}. The baryon number corresponds to the
  topological charge.
\item
  \textbf{Equivalence between the quantum sine-Gordon model and the
  massive Thirring model}: Coleman {[}20{]} and Mandelstam {[}21{]}
  showed that the quantum theory of the sine-Gordon model is equivalent
  to the massive Thirring model (a fermionic field theory). The
  \textbf{boson-fermion correspondence}, in which bosonic solitons
  correspond to fermions, is one of the most remarkable results of
  1+1-dimensional field theory.
\end{itemize}

\subsubsection{6.3 Differential Geometry}\label{differential-geometry}

\begin{itemize}
\tightlist
\item
  \textbf{Pseudospheres}: The compatibility conditions between the first
  fundamental form and the second fundamental form of surfaces with
  constant negative curvature (\(K = -1\)) reduce to the sine-Gordon
  equation {[}5, 6{]}. Kink solutions correspond to geodesics on
  pseudospheres.
\end{itemize}

\begin{center}\rule{0.5\linewidth}{0.5pt}\end{center}

\subsection{7. Implications for Physical
Interpretation}\label{implications-for-physical-interpretation}

\subsubsection{7.1 Summary of Particle-like Properties of
Solitons}\label{summary-of-particle-like-properties-of-solitons}

The soliton solutions of the sine-Gordon equation naturally possess the
following particle-like properties:

{\def\LTcaptype{none} % do not increment counter
\begin{longtable}[]{@{}
  >{\centering\arraybackslash}p{(\linewidth - 2\tabcolsep) * \real{0.4194}}
  >{\centering\arraybackslash}p{(\linewidth - 2\tabcolsep) * \real{0.5806}}@{}}
\toprule\noalign{}
\begin{minipage}[b]{\linewidth}\centering
Soliton Property
\end{minipage} & \begin{minipage}[b]{\linewidth}\centering
Correspondence to Particles
\end{minipage} \\
\midrule\noalign{}
\endhead
\bottomrule\noalign{}
\endlastfoot
Winding number \(N = \pm 1\) & Particle / anti-particle \\
Conservation of winding number & Conservation of particle number \\
Shape preservation during collisions & Identity of particles \\
Phase shift (position change after collision) & Scattering through
interactions \\
Lorentz covariance & Special relativistic structure \\
Breather (\(N = 0\) bound state) & Meson-like bound state \\
Periodicity of the potential (\(2\pi\)) & Discreteness of phase \\
Boson-fermion correspondence {[}20, 21{]} & Spin-statistics \\
\end{longtable}
}

\subsubsection{7.2 Significance of the Rectangular Wave
Description}\label{significance-of-the-rectangular-wave-description}

As shown in §8--§9, all of these particle-like properties are preserved
in the rectangular wave limit \(\sigma \to 0\). This means that the
essential information of solitons is not contained in the smooth shape
of the continuous field, but is concentrated in the topological and
scattering discrete structure.

\begin{center}\rule{0.5\linewidth}{0.5pt}\end{center}

\subsection{8. Finite Parameter Representation of
Solitons}\label{finite-parameter-representation-of-solitons}

\subsubsection{8.1 Motivation: From Continuous Fields to Finite
Parameters}\label{motivation-from-continuous-fields-to-finite-parameters}

Solutions of the sine-Gordon equation are defined as continuous fields
\(\phi(x, t)\), but due to their localization, soliton solutions are
\textbf{completely characterized by a finite number of parameters}. This
is not coincidental. Solitons are \textbf{finite-mode solutions} in the
mode decomposition of nonlinear fields, solutions restricted to
finite-dimensional invariant manifolds (moduli spaces) within the
infinite-dimensional field theory {[}10, 22{]}.

In this section, we provide a formulation that describes the soliton
solutions of the sine-Gordon equation with a finite number of
parameters, and explicitly show how solutions of the continuous wave
equation are reduced to a small number of discrete pieces of
information.

\subsubsection{8.2 Five-Parameter Representation of the
Kink}\label{five-parameter-representation-of-the-kink}

We reproduce the kink solution from §2.1:

\[\phi_K(x, t) = 4\arctan\left(\exp\left(\frac{x - x_0 - vt}{\ell}\right)\right)\]

Here \(\ell = \sqrt{1 - v^2}\) is the width parameter of the kink. This
solution is \textbf{completely} determined by the following five
parameters.

\textbf{Definition 8.1 (Soliton parameters).} The five parameters
characterizing a 1+1-dimensional sine-Gordon kink are defined as
follows:

{\def\LTcaptype{none} % do not increment counter
\begin{longtable}[]{@{}
  >{\centering\arraybackslash}p{(\linewidth - 6\tabcolsep) * \real{0.2500}}
  >{\centering\arraybackslash}p{(\linewidth - 6\tabcolsep) * \real{0.2500}}
  >{\centering\arraybackslash}p{(\linewidth - 6\tabcolsep) * \real{0.2500}}
  >{\centering\arraybackslash}p{(\linewidth - 6\tabcolsep) * \real{0.2500}}@{}}
\toprule\noalign{}
\begin{minipage}[b]{\linewidth}\centering
Symbol
\end{minipage} & \begin{minipage}[b]{\linewidth}\centering
Name
\end{minipage} & \begin{minipage}[b]{\linewidth}\centering
Definition
\end{minipage} & \begin{minipage}[b]{\linewidth}\centering
Range
\end{minipage} \\
\midrule\noalign{}
\endhead
\bottomrule\noalign{}
\endlastfoot
\(N\) & Winding number & Topological charge \(\pm 1\) &
\(\{-1, +1\}\) \\
\(A\) & Maximum wave height & Total change in \(\phi\) \(= 2\pi |N|\) &
\(2\pi\) (fixed) \\
\(\theta_0\) & Center phase & Spatial phase of soliton center
\(= 2\pi x_0 / L\) & \([0, 2\pi)\) \\
\(\alpha\) & Velocity phase & Phase angle representation of velocity
\(v = \sin\alpha\) & \((-\pi/2, \pi/2)\) \\
\(\sigma\) & Spread parameter & Phase angle representation of width
\(\ell = \cos\alpha\) & \((0, 1]\) \\
\end{longtable}
}

\textbf{Proposition 8.1 (Completeness of parameters).} A single kink
(anti-kink) solution of the sine-Gordon equation is uniquely determined
by the above five parameters \((N, A, \theta_0, \alpha, \sigma)\).
Conversely, the above five parameters can be uniquely read from any kink
solution.

\emph{Proof.} The general form of the kink solution is
\(\phi = 4\arctan(\exp(\pm(x - x_0 - vt)/\sqrt{1-v^2}))\), where the
sign \(\pm\) determines \(N\), the displacement \(x_0\) determines
\(\theta_0\), and the velocity \(v\) determines \(\alpha\). \(A = 2\pi\)
is fixed from the winding number \(|N| = 1\), and
\(\sigma = \cos\alpha\) is determined subordinately from \(\alpha\). The
free parameters are the three quantities \((N, \theta_0, \alpha)\), and
\(A\) and \(\sigma\) are derived from these. \(\square\)

\textbf{Remark.} The free parameters are essentially three
(\(N, \theta_0, \alpha\)), but the five-parameter representation is
adopted for clarity of physical interpretation. \(A\) is a dependent
parameter determined from the winding number, and \(\sigma\) from the
velocity phase.

\subsubsection{8.3 Relationship Between Velocity Phase and Spread
Parameter}\label{relationship-between-velocity-phase-and-spread-parameter}

The following identity holds between the velocity phase \(\alpha\) and
the spread parameter \(\sigma\):

\[\sin^2\alpha + \sigma^2 = v^2 + \ell^2 = 1\]

That is, \((\sin\alpha, \sigma) = (v, \ell)\) is a point on the unit
circle.

\textbf{Physical interpretation:}

\begin{itemize}
\tightlist
\item
  At \(\alpha = 0\) (rest), \(\sigma = 1\): the soliton has maximum
  width
\item
  As \(\alpha \to \pm\pi/2\) (speed of light), \(\sigma \to 0\): the
  width of the soliton contracts to zero (Lorentz contraction)
\item
  The faster the motion, the narrower the spatial extent --
  \textbf{motion and shape are unified by a single phase angle
  \(\alpha\)}
\end{itemize}

\subsubsection{8.4 Six-Parameter Representation of the
Breather}\label{six-parameter-representation-of-the-breather}

The breather solution of §2.4 requires the internal oscillation
frequency \(\omega\) in addition to the five kink parameters.

\textbf{Definition 8.2 (Breather parameters).}

{\def\LTcaptype{none} % do not increment counter
\begin{longtable}[]{@{}
  >{\centering\arraybackslash}p{(\linewidth - 6\tabcolsep) * \real{0.2500}}
  >{\centering\arraybackslash}p{(\linewidth - 6\tabcolsep) * \real{0.2500}}
  >{\centering\arraybackslash}p{(\linewidth - 6\tabcolsep) * \real{0.2500}}
  >{\centering\arraybackslash}p{(\linewidth - 6\tabcolsep) * \real{0.2500}}@{}}
\toprule\noalign{}
\begin{minipage}[b]{\linewidth}\centering
Symbol
\end{minipage} & \begin{minipage}[b]{\linewidth}\centering
Name
\end{minipage} & \begin{minipage}[b]{\linewidth}\centering
Definition
\end{minipage} & \begin{minipage}[b]{\linewidth}\centering
Range
\end{minipage} \\
\midrule\noalign{}
\endhead
\bottomrule\noalign{}
\endlastfoot
\(N\) & Winding number & \(0\) (kink-anti-kink bound state) & \(\{0\}\)
(fixed) \\
\(A\) & Maximum wave height & \(4\arctan(\sqrt{1-\omega^2}/\omega)\) &
\((0, \pi)\) \\
\(\theta_0\) & Center phase & Spatial phase of breather center &
\([0, 2\pi)\) \\
\(\alpha\) & Velocity phase & Overall motion \(v = \sin\alpha\) &
\((-\pi/2, \pi/2)\) \\
\(\sigma\) & Spread parameter & \(1/\sqrt{1-\omega^2}\) &
\([1, \infty)\) \\
\(\omega\) & Internal phase velocity & Angular frequency of internal
oscillation & \((0, 1)\) \\
\end{longtable}
}

The breather has \(N = 0\) but possesses the additional degree of
freedom of internal oscillation \(\omega\).

\subsubsection{8.5 Parameter Representation of Multi-Soliton
States}\label{parameter-representation-of-multi-soliton-states}

A multi-body state consisting of \(n\) solitons is represented as a set
of parameters for each soliton {[}10, 12{]}.

\textbf{Definition 8.3 (\(n\)-soliton state).} A state consisting of
\(n\) solitons is defined by the following parameter set:

\[\mathcal{S}_n = \{(N_i, A_i, \theta_{0,i}, \alpha_i, \sigma_i)\}_{i=1}^{n}\]

\textbf{Proposition 8.2 (Asymptotic completeness).} A field
configuration \(\phi(x, t)\) consisting of \(n\) sufficiently separated
solitons is asymptotically completely determined by the parameter set
\(\mathcal{S}_n\). In the region where overlap between solitons is
negligible,

\[\phi(x, t) \approx \sum_{i=1}^{n} \phi_i(x, t; N_i, \theta_{0,i}, \alpha_i)\]

holds, where \(\phi_i\) is the single-soliton solution for each soliton.

\emph{Proof.} By the inverse scattering method {[}12{]}, the
\(n\)-soliton solution is completely determined by \(n\) scattering data
(eigenvalues and norming constants). In the asymptotic region, the
scattering data and the above parameters are in one-to-one
correspondence. \(\square\)

\subsubsection{8.6 Description of Collisions in Parameter
Space}\label{description-of-collisions-in-parameter-space}

The soliton collisions of §3 can be described as mappings on parameter
space.

\textbf{Proposition 8.3 (Collision mapping).} The collision of two
solitons \(\mathcal{S}_1 = (N_1, \theta_{0,1}, \alpha_1)\) and
\(\mathcal{S}_2 = (N_2, \theta_{0,2}, \alpha_2)\) is described as the
following parameter transformation:

\[\mathcal{S}_1' = (N_1,\; \theta_{0,1} + \Delta\theta_1,\; \alpha_1)\]
\[\mathcal{S}_2' = (N_2,\; \theta_{0,2} + \Delta\theta_2,\; \alpha_2)\]

That is:

\begin{enumerate}
\def\labelenumi{\arabic{enumi}.}
\tightlist
\item
  \textbf{Winding numbers \(N_i\) are invariant} -- the type of particle
  does not change
\item
  \textbf{Velocity phases \(\alpha_i\) are invariant} -- velocity
  (momentum) is conserved
\item
  \textbf{Only the center phases \(\theta_{0,i}\) change} -- they are
  displaced by the phase shift \(\Delta\theta_i\)
\end{enumerate}

\textbf{Corollary 8.1.} All information about a soliton collision is
concentrated in the phase shifts \(\Delta\theta_i\). The phase shift is
computable from the incident parameters {[}12{]} and takes the form

\[\Delta\theta_1 = -\Delta\theta_2 = f(N_1, N_2, \alpha_1 - \alpha_2)\]

That is, the collision depends only on the relative velocity phase
\(\alpha_1 - \alpha_2\).

\subsubsection{8.7 Reduction from Continuous Fields to Parameter
Representation}\label{reduction-from-continuous-fields-to-parameter-representation}

Summarizing the above, the solution space of the sine-Gordon equation is
stratified as follows:

{\def\LTcaptype{none} % do not increment counter
\begin{longtable}[]{@{}
  >{\centering\arraybackslash}p{(\linewidth - 4\tabcolsep) * \real{0.4286}}
  >{\centering\arraybackslash}p{(\linewidth - 4\tabcolsep) * \real{0.2857}}
  >{\centering\arraybackslash}p{(\linewidth - 4\tabcolsep) * \real{0.2857}}@{}}
\toprule\noalign{}
\begin{minipage}[b]{\linewidth}\centering
Description Level
\end{minipage} & \begin{minipage}[b]{\linewidth}\centering
Degrees of Freedom
\end{minipage} & \begin{minipage}[b]{\linewidth}\centering
Information Content
\end{minipage} \\
\midrule\noalign{}
\endhead
\bottomrule\noalign{}
\endlastfoot
Continuous field \(\phi(x, t)\) & Infinite-dimensional & Real value at
each spacetime point \\
\(n\)-soliton solution & \(3n\)-dimensional & Each soliton
\((N_i, \theta_{0,i}, \alpha_i)\) \\
Change after collision & \(n\)-dimensional & Only phase shifts
\(\Delta\theta_i\) \\
\end{longtable}
}

\textbf{Proposition 8.4 (Finite-dimensionality of the soliton sector).}
Within the solution space of the sine-Gordon equation, the sector
containing \(n\) solitons is described as a \(3n\)-dimensional
finite-dimensional manifold (moduli space) \(\mathcal{M}_n\) {[}10,
22{]}. The infinite degrees of freedom of the continuous field are
\textbf{rigorously} reduced to a finite number of parameters in the
soliton sector.

\emph{Proof.} In the inverse scattering method, the \(n\)-soliton sector
is characterized by \(n\) discrete eigenvalues \(\zeta_i\) and their
associated norming constants \(c_i\) {[}12, 13{]}. Each eigenvalue
determines the winding number \(N_i\) and velocity \(v_i\) (=
\(\sin\alpha_i\)), and the norming constants determine the initial
positions \(x_{0,i}\) (= \(\theta_{0,i}\)). Excluding the radiation
component (continuous spectrum), the degrees of freedom of the soliton
sector are exactly \(3n\). \(\square\)

This reduction is not an approximation. In the soliton sector, the
infinite-dimensional partial differential equation \textbf{rigorously}
reduces to a finite-dimensional system of ordinary differential
equations.

\subsubsection{8.8 Relationship with Rectangular
Waves}\label{relationship-with-rectangular-waves}

Let us re-examine the shape of the kink solution:

\[\phi_K(x) = 4\arctan(e^{x/\ell}) = 2\pi \cdot \frac{1}{1 + e^{-x/\ell}}\]

The right-hand side is a logistic function (sigmoid function), which
converges to a \textbf{step function} (half-period of a rectangular
wave) in the limit \(\ell \to 0\):

\[\lim_{\ell \to 0} \phi_K(x) = \begin{cases} 0 & (x < 0) \\ 2\pi & (x > 0) \end{cases}\]

\textbf{Proposition 8.5 (Rectangular wave limit).} In the limit of the
spread parameter \(\sigma = \ell \to 0\) (i.e., \(\alpha \to \pm\pi/2\),
the speed-of-light limit), the kink solution converges to a rectangular
wave (step function). Conversely, a kink with finite width
\(\sigma > 0\) can be understood as a \textbf{smoothing} of the
rectangular wave.

\emph{Proof.} \(\phi_K(x) = 2\pi/(1 + e^{-x/\ell})\) is a logistic
function, and converges pointwise to the Heaviside step function
\(2\pi \cdot H(x)\) as \(\ell \to 0\). \(\square\)

\textbf{Corollary 8.2.} A soliton can be interpreted as a finite-mode
approximation of a rectangular wave. The continuous \(\tanh\)-type shape
is merely a smoothed version of a discrete step function. The spread
parameter \(\sigma\) controls the degree of smoothing.

From this viewpoint, the parameter representation of solitons is
reinterpreted as follows:

{\def\LTcaptype{none} % do not increment counter
\begin{longtable}[]{@{}
  >{\centering\arraybackslash}p{(\linewidth - 4\tabcolsep) * \real{0.2500}}
  >{\centering\arraybackslash}p{(\linewidth - 4\tabcolsep) * \real{0.2917}}
  >{\centering\arraybackslash}p{(\linewidth - 4\tabcolsep) * \real{0.4583}}@{}}
\toprule\noalign{}
\begin{minipage}[b]{\linewidth}\centering
Parameter
\end{minipage} & \begin{minipage}[b]{\linewidth}\centering
Meaning in Continuous Field
\end{minipage} & \begin{minipage}[b]{\linewidth}\centering
Meaning in Rectangular Wave (Discrete)
\end{minipage} \\
\midrule\noalign{}
\endhead
\bottomrule\noalign{}
\endlastfoot
\(N = \pm 1\) & Winding number & Direction of step
(ascending/descending) \\
\(A = 2\pi\) & Total phase change & Height of step \\
\(\theta_0\) & Position of soliton center & Edge position of step \\
\(\alpha\) & Phase angle of velocity & Direction of edge movement \\
\(\sigma\) & Width of smoothing & \textbf{Equals rectangular wave at
\(\sigma \to 0\)} \\
\end{longtable}
}

\begin{center}\rule{0.5\linewidth}{0.5pt}\end{center}

\subsection{9. Preservation of Properties in the Rectangular Wave
Limit}\label{preservation-of-properties-in-the-rectangular-wave-limit}

\subsubsection{9.1 Purpose}\label{purpose}

The properties of the sine-Gordon equation described in §1--§7 have been
proved for the continuous field \(\phi(x, t)\). As shown in §8.8, the
soliton solution converges to a rectangular wave (step function) in the
limit of the spread parameter \(\sigma \to 0\).

In this section, for each property from §1--§7, we systematically verify
that (a) it holds for soliton solutions (which is trivial since they are
exact solutions), and (b) it is preserved in the rectangular wave limit
\(\sigma \to 0\).

If all properties are preserved in the rectangular wave limit, then the
rectangular wave description is not a ``convenient approximation'' but
rather a \textbf{rigorous limit that completely inherits all structural
properties} of the soliton.

\subsubsection{9.2 Framework for
Verification}\label{framework-for-verification}

Each property is verified in the following three stages:

\begin{enumerate}
\def\labelenumi{\arabic{enumi}.}
\tightlist
\item
  \textbf{Validity for soliton solutions}: Does it satisfy the
  sine-Gordon equation as an exact solution?
\item
  \textbf{\(\sigma\)-dependence analysis}: Does the property depend on
  the spread parameter \(\sigma\)?
\item
  \textbf{Existence of the rectangular wave limit}: Is the property
  preserved or does it diverge as \(\sigma \to 0\)?
\end{enumerate}

\subsubsection{9.3 Definition and Conservation of Winding Number (§2.3 →
Rectangular
Wave)}\label{definition-and-conservation-of-winding-number-2.3-rectangular-wave}

\textbf{Property from §2.3:} The winding number
\(N = [\phi(+\infty) - \phi(-\infty)] / 2\pi\) is conserved under time
evolution.

\textbf{Validity for soliton solutions:} The kink solution gives
\(\phi(-\infty) = 0\), \(\phi(+\infty) = 2\pi\), so \(N = +1\). Since
this is an exact solution of the sine-Gordon equation, the conservation
of \(N\) is trivial. \(\checkmark\)

\textbf{\(\sigma\)-dependence:} The definition of winding number
\(N = [\phi(+\infty) - \phi(-\infty)] / 2\pi\) depends \textbf{only on
the boundary values} and does not depend at all on the internal shape of
\(\phi\) (hence on \(\sigma\)).

\textbf{Proposition 9.1 (\(\sigma\)-independence of the winding
number).} The winding number \(N\) does not depend on the spread
parameter \(\sigma\). Even in the rectangular wave limit
\(\sigma \to 0\), \(N\) retains the same integer value.

\emph{Proof.} For any \(\sigma > 0\), \(\phi_K^\sigma(\pm\infty)\) are
fixed at the vacuum values \(0, 2\pi\). In the limit \(\sigma \to 0\),
\(\phi_K^\sigma(x) \to 2\pi \cdot H(x)\) converges pointwise, but since
\(H(-\infty) = 0\), \(H(+\infty) = 1\), we have
\(N = (2\pi \cdot 1 - 2\pi \cdot 0)/2\pi = 1\), which is invariant.
\(\square\)

\textbf{Winding number for rectangular waves:} For the rectangular wave
\(\phi(x) = 2\pi H(x - x_0)\), \(N = +1\). If the step is ascending,
\(N = +1\); if descending, \(N = -1\). The winding number is precisely
the sign of the step itself. \(\checkmark\)

\subsubsection{9.4 Conservation Law for Winding Number (Proposition 2.1
→ Rectangular
Wave)}\label{conservation-law-for-winding-number-proposition-2.1-rectangular-wave}

\textbf{Property from §2.3:} \(N\) is conserved under time evolution.
Creation and annihilation are possible only in kink-anti-kink pairs.

\textbf{\(\sigma\)-dependence:} The proof of the conservation law (§2.3)
relies only on the fact that ``the boundary values are fixed at vacuum
values and cannot change by integer multiples of \(2\pi\) through
continuous time evolution.'' This argument does not depend on
\(\sigma\).

\textbf{Conservation for rectangular waves:} The boundary values of the
rectangular wave are \(0\) and \(2\pi\), taking only discrete values.
Discrete values cannot change under continuous deformation. Therefore,
the conservation of winding number for rectangular waves is
\textbf{rather more obvious} than for the continuous field case.

\textbf{Proposition 9.2.} The conservation law for the winding number
holds rigorously in the rectangular wave limit. Furthermore, since
\(\phi\) itself takes only discrete values (\(0\) or \(2\pi\)) in the
rectangular wave case, conservation is guaranteed not only topologically
but also by the \textbf{discreteness of values}. \(\checkmark\)

\subsubsection{9.5 Existence of Kink and Anti-kink Solutions (§2.1, §2.2
→ Rectangular
Wave)}\label{existence-of-kink-and-anti-kink-solutions-2.1-2.2-rectangular-wave}

\textbf{Property from §2.1--2.2:} Localized solutions with winding
number \(\pm 1\) exist and propagate with velocity \(v\).

\textbf{Validity for soliton solutions:} The kink
\(\phi_K = 4\arctan(\exp((x-vt)/\sigma))\) is an exact solution.
\(\checkmark\)

\textbf{Rectangular wave limit:} In the limit \(\sigma \to 0\),
\(\phi_K \to 2\pi \cdot H(x - vt)\). This is a step function moving with
velocity \(v\), preserving localization and propagation properties.

\textbf{Proposition 9.3 (Existence of rectangular wave kink).} The
rectangular wave kink \(\phi(x, t) = 2\pi \cdot H(x - vt)\) has the
following properties: 1. Winding number \(N = +1\) 2. Moves with
velocity \(v\) 3. Localized at a single point in space (\(x = vt\)) 4.
Holds as a \textbf{distributional solution} of the sine-Gordon equation

\emph{Proof.} (1)--(3) are obvious from the definition. For (4): the
distributional derivatives of \(\phi = 2\pi H(x - vt)\) are
\(\phi_x = 2\pi\delta(x-vt)\), \(\phi_{xx} = 2\pi\delta'(x-vt)\),
\(\phi_{tt} = 2\pi v^2\delta'(x-vt)\), and
\(\sin(\phi) = \sin(2\pi H(x-vt)) = 0\) (for \(x \neq vt\), \(\phi = 0\)
or \(2\pi\), both giving \(\sin = 0\)). The wave equation part
\(\phi_{tt} - \phi_{xx} = 2\pi(v^2-1)\delta'(x-vt)\) vanishes only when
\(v = \pm 1\) (speed of light). For \(v \neq 1\), the rectangular wave
is not a classical solution, and smoothing with finite width
\(\sigma > 0\) is necessary. \(\square\)

\textbf{Remark.} The fact that the rectangular wave is not an exact
solution for \(v \neq 1\) is consistent with the relation
\(\sigma = \cos\alpha = \sqrt{1-v^2}\) from §8.3. If \(v < 1\), then
\(\sigma > 0\), and the soliton has finite width. The rectangular wave
(\(\sigma = 0\)) corresponds only to \(v = 1\) (speed of light). The
fact that particles with velocities below the speed of light have a
finite ``spread'' is qualitatively consistent with the spread of wave
packets in quantum mechanics.

\subsubsection{9.6 Breather Solution (§2.4 → Rectangular
Wave)}\label{breather-solution-2.4-rectangular-wave}

\textbf{Property from §2.4:} A breather exists as a kink-anti-kink bound
state.

\textbf{Rectangular wave limit:} The breather is an oscillating solution
with winding number \(N = 0\). In the rectangular wave limit, this
corresponds to a \textbf{rectangular pulse} oscillating between \(0\)
and \(2\pi\).

\textbf{Proposition 9.4.} The rectangular wave limit of the breather is
the oscillation of a spatially localized rectangular pulse. It
periodically transitions between two states: presence (\(\phi = 2\pi\))
and absence (\(\phi = 0\)) of the pulse. The mass
\(M_B = 16\sqrt{1-\omega^2}\) depends on \(\omega\) and does not
directly depend on \(\sigma\). \(\checkmark\)

\subsubsection{9.7 Soliton Collisions and Phase Shift (§3 → Rectangular
Wave)}\label{soliton-collisions-and-phase-shift-3-rectangular-wave}

\textbf{Property from §3:} After collision: (i) shape preservation, (ii)
phase shift, (iii) winding number conservation.

This is the most important property, and each item is verified
individually.

\textbf{(i) Shape preservation:}

\textbf{Soliton solution:} The soliton after collision has the same
shape (same \(\sigma\)) as before the collision {[}12{]}. \(\checkmark\)

\textbf{Rectangular wave limit:} The rectangular wave has
\(\sigma = 0\), and the only shape parameter is ``step height
\(2\pi\).'' If \(\sigma = 0\) is preserved before and after the
collision, shape preservation is automatically satisfied. Since
\(\sigma\) is a conserved quantity in soliton collisions (Proposition
8.3: \(\alpha_i\) invariant \(\Rightarrow\) \(\sigma_i = \cos\alpha_i\)
invariant), \(\sigma = 0\) is also conserved. \(\checkmark\)

\textbf{(ii) Phase shift:}

\textbf{Soliton solution:} The phase shift after collision is
\(\Delta\theta = f(N_1, N_2, \alpha_1 - \alpha_2)\) {[}12{]}.
\(\checkmark\)

\textbf{Rectangular wave limit:} The phase shift formula depends only on
\(\alpha_i\) (velocity phase) and does not depend on \(\sigma_i\)
(spread) (see §8.6, Corollary 8.1). Therefore, even at \(\sigma \to 0\),
the value of the phase shift does not change.

\textbf{Proposition 9.5 (\(\sigma\)-independence of phase shift).} The
phase shift \(\Delta\theta\) of a soliton collision does not depend on
the spread parameter \(\sigma_i\) of each soliton. Even in the
rectangular wave limit \(\sigma \to 0\), the phase shift takes the same
value as in the soliton case.

\emph{Proof.} In the inverse scattering method, the phase shift is
computed only from the discrete eigenvalues \(\zeta_i\) of the
scattering data {[}12, 13{]}. The \(\zeta_i\) are determined by the
soliton velocity \(v_i = \sin\alpha_i\) and winding number \(N_i\), and
are independent of the norming constants (position information) and
radiation component. The soliton width \(\sigma_i\) is related to the
norming constants but does not contribute to the phase shift.
\(\square\)

\textbf{(iii) Winding number conservation:} Already shown in §9.4.
\(\checkmark\)

\subsubsection{9.8 Complete Integrability and Conserved Quantities (§4 →
Rectangular
Wave)}\label{complete-integrability-and-conserved-quantities-4-rectangular-wave}

\textbf{Property from §4:} Conservation of energy, momentum, and winding
number.

\textbf{Energy:}

The energy of a soliton is
\(E = 8/\sigma = 8/\cos\alpha = 8/\sqrt{1-v^2}\). As \(\sigma \to 0\),
the energy diverges.

\textbf{Proposition 9.6 (\(\sigma\)-dependence of energy).} The energy
of a soliton depends on \(\sigma\), and diverges as \(E \to \infty\)
when \(\sigma \to 0\). This means that the rectangular wave limit has
infinite energy.

\textbf{Proposition 9.6a (Limits of applicability of this theory).} The
sine-Gordon soliton theory is based on the following implicit
assumptions:

\begin{enumerate}
\def\labelenumi{\arabic{enumi}.}
\tightlist
\item
  \textbf{Fixed origin}: The initial value of the field is fixed at the
  vacuum value \(\phi(-\infty) = 0\) (§2.1).
\item
  \textbf{Periodicity of the potential}: \(V(\phi) = 1-\cos\phi\) has
  period \(2\pi\), and the target space is effectively \(S^1\). The size
  of one ``circuit'' of the potential (\(2\pi\)) is fixed in the theory
  and has no freedom for modification.
\item
  \textbf{Forcing the terminus to vacuum values}: The requirement of
  finite energy on the infinite space \((-\infty, +\infty)\) forces
  \(\phi(+\infty) = 2n\pi\) (\(n \in \mathbb{Z}\)). This is because
  \(V(2n\pi) = 0\) (potential minimum), while \(\phi = (2n+1)\pi\) gives
  \(V = 2\) (potential maximum), so energy diverges on infinite space,
  and such configurations are excluded.
\item
  \textbf{Exclusion of fractional winding numbers}: As a consequence of
  the above, the winding number \(N\) is always an integer. Half-integer
  winding numbers (\(\phi(+\infty) = (2n+1)\pi\)) are automatically
  excluded by the finite energy requirement.
\item
  \textbf{Absence of boundary reflection}: Since the space is infinite
  or closed (\(S^1\)), there are no physical boundaries, and boundary
  reflection of solitons does not occur.
\item
  \textbf{Perfectly elastic collisions and absence of interactions}:
  This theory is a completely integrable system (§4) possessing
  infinitely many conserved quantities. As a consequence of this
  excessive constraint, soliton collisions are described as perfectly
  elastic collisions, where the velocity, shape, and energy of each
  soliton are \textbf{individually} conserved before and after the
  collision (§3.3). To avoid violating momentum conservation, only the
  position of each soliton (phase shift \(\Delta x\)) changes, and
  exchange of velocities, deformation of shapes, and dissipation of
  energy never occur. This effectively means that solitons behave as
  \textbf{free particles that do not interact} with each other. For
  identical particles (e.g., kink-kink collision), ``whether they passed
  through or repelled'' is in principle indistinguishable.
\end{enumerate}

\emph{Consequence.} This theory is consistently applicable only in the
following two cases:

\begin{itemize}
\tightlist
\item
  \textbf{Infinite space} \((-\infty, +\infty)\): The endpoints are at
  infinity and are never reached. The energy is \(E = 8\) per kink (at
  rest) and finite.
\item
  \textbf{Closed curve} \(S^1\) (circumference \(L\)): Periodic boundary
  conditions \(\phi(x+L) = \phi(x) + 2\pi N\). No boundaries exist.
\end{itemize}

An \textbf{open string} (Dirichlet / Neumann boundary conditions on a
finite interval \([0, L]\)) lies outside the scope of this theory. The
reason is that a finite interval \([0, L]\) is a contractible space
(\(\pi_1([0,L]) = 0\)), so the winding number is not defined as a
topological invariant, and the topological stability of solitons is not
guaranteed. Furthermore, on a finite interval, \(\phi(L) = (2n+1)\pi\)
(half-integer winding number, potential maximum) is energetically
allowed, so the restriction to integer winding numbers does not hold.

\textbf{Self-cancellation structure of energy:} The energy of each
\(2\pi\) rotation (one kink) is self-contained at \(E = 8\) (at rest).
This is because the potential \(V(\phi) = 1-\cos\phi\) returns from
\(0 \to 0\) over one period. A kink-anti-kink pair (\(4\pi\) rotation,
winding number \(N = 0\)) can return its energy to zero through pair
annihilation. The only observable permanent energy is the contribution
\(8|N|\) from the remaining kinks that cannot undergo pair annihilation.

\textbf{Momentum and winding number:} The momentum \(P = -Ev\) has a
similar structure. The winding number is \(\sigma\)-independent as shown
in §9.3. \(\checkmark\)

\subsubsection{9.9 Extension to Higher Dimensions (§5 → Rectangular
Wave)}\label{extension-to-higher-dimensions-5-rectangular-wave}

\textbf{Property from §5:} By Derrick's theorem, stable localized
solutions do not exist in infinite space of 3+1 dimensions.

\textbf{Effect on rectangular wave limit:} The premise of Derrick's
theorem is the \textbf{existence of finite-energy solutions in infinite
space} {[}17{]}. On a closed curve \(S^1\) (circumference \(L\)), the
premises of the theorem break down for the following two reasons:

\begin{enumerate}
\def\labelenumi{\arabic{enumi}.}
\tightlist
\item
  \textbf{Space is closed}: The range of \(\lambda\) in the scale
  transformation \(\phi(x) \to \phi(\lambda x)\) is restricted
\item
  \textbf{Topological protection}: The winding number
  \(N \in \pi_1(S^1) = \mathbb{Z}\) is protected as a topological
  invariant
\end{enumerate}

\textbf{Proposition 9.7.} On a closed curve \(S^1\) (circumference
\(L\)), localized solutions of any width, including rectangular waves
(\(\sigma = 0\)), are topologically stable. Derrick's theorem does not
apply to closed spaces. \(\checkmark\)

\subsubsection{9.10 Summary of Verification
Results}\label{summary-of-verification-results}

{\def\LTcaptype{none} % do not increment counter
\begin{longtable}[]{@{}
  >{\centering\arraybackslash}p{(\linewidth - 10\tabcolsep) * \real{0.1132}}
  >{\centering\arraybackslash}p{(\linewidth - 10\tabcolsep) * \real{0.0755}}
  >{\centering\arraybackslash}p{(\linewidth - 10\tabcolsep) * \real{0.1698}}
  >{\centering\arraybackslash}p{(\linewidth - 10\tabcolsep) * \real{0.3019}}
  >{\centering\arraybackslash}p{(\linewidth - 10\tabcolsep) * \real{0.2264}}
  >{\centering\arraybackslash}p{(\linewidth - 10\tabcolsep) * \real{0.1132}}@{}}
\toprule\noalign{}
\begin{minipage}[b]{\linewidth}\centering
Property
\end{minipage} & \begin{minipage}[b]{\linewidth}\centering
Section
\end{minipage} & \begin{minipage}[b]{\linewidth}\centering
Soliton
\end{minipage} & \begin{minipage}[b]{\linewidth}\centering
\(\sigma\)-dependence
\end{minipage} & \begin{minipage}[b]{\linewidth}\centering
Rectangular Wave Limit
\end{minipage} & \begin{minipage}[b]{\linewidth}\centering
Remarks
\end{minipage} \\
\midrule\noalign{}
\endhead
\bottomrule\noalign{}
\endlastfoot
Definition of winding number & §2.3 & \(\checkmark\) &
\textbf{Independent} & \(\checkmark\) & Boundary values only \\
Conservation of winding number & §2.3 & \(\checkmark\) &
\textbf{Independent} & \(\checkmark\) & Topological \\
Existence of kink/anti-kink & §2.1--2 & \(\checkmark\) & Shape only &
\(\checkmark\) & Distributional solution \\
Existence of breather & §2.4 & \(\checkmark\) & Shape only &
\(\checkmark\) & Rectangular pulse \\
Shape preservation after collision & §3.1--2 & \(\checkmark\) &
\textbf{Independent} & \(\checkmark\) & \(\sigma\) conserved \\
Phase shift after collision & §3.3 & \(\checkmark\) &
\textbf{Independent} & \(\checkmark\) & \(\alpha\) only \\
Winding number conservation (collision) & §3.3 & \(\checkmark\) &
\textbf{Independent} & \(\checkmark\) & Topological \\
Energy conservation & §4.2 & \(\checkmark\) & Dependent &
\(\checkmark\)* & *Self-contained per \(2\pi\) rotation \\
Momentum conservation & §4.2 & \(\checkmark\) & Dependent &
\(\checkmark\)* & *Same as above \\
Derrick avoidance & §5.1--2 & -- & -- & \(\checkmark\) & Closed space
\(S^1\) \\
Parameter representation \((N, \theta_0, \alpha)\) & §8 & \(\checkmark\)
& \textbf{Independent} & \(\checkmark\) & \(\sigma\) is subordinate \\
\end{longtable}
}

\textbf{Theorem 9.1 (Structural preservation of the rectangular wave
limit).} The following structural properties of soliton solutions of the
sine-Gordon equation are all preserved in the rectangular wave limit
\(\sigma \to 0\):

\begin{enumerate}
\def\labelenumi{\arabic{enumi}.}
\tightlist
\item
  \textbf{Definition and conservation of winding number} (topological
  property: \(\sigma\)-independent)
\item
  \textbf{Shape preservation after collision and phase shift} (property
  of scattering data: \(\sigma\)-independent)
\item
  \textbf{Finite parameter representation} \((N, \theta_0, \alpha)\)
  (structure of moduli space: \(\sigma\)-independent)
\item
  \textbf{Finiteness of energy and momentum} (due to the
  self-cancellation structure of each \(2\pi\) rotation; Proposition
  9.6a)
\end{enumerate}

\emph{Proof.} (1) follows from Propositions 9.1--9.2, (2) from
Proposition 9.5, (3) from Proposition 8.4 and the subordination
\(\sigma = \cos\alpha\), and (4) from the energy self-cancellation
structure of Proposition 9.6a. \(\square\)

\textbf{Corollary 9.1.} The rectangular wave description is not a
``coarse approximation'' of the soliton. It is a \textbf{rigorous limit
that completely inherits} the topological and scattering structure of
the soliton; what is lost is only the spread parameter \(\sigma\)
(details of the internal shape).

\textbf{Corollary 9.2 (Limits of applicability).} The conditions for the
validity of this theorem depend on the implicit assumptions of the
theory made explicit in Proposition 9.6a -- the fixed origin
\(\phi = 0\), the forcing of the terminus to vacuum values (\(2n\pi\)),
and integer winding numbers. These assumptions hold only on infinite
space \((-\infty,+\infty)\) or on a closed curve \(S^1\). Extension to
open strings (boundary conditions on a finite interval) requires the
admission of half-integer winding numbers and a theory of soliton
reflection at boundaries (Ghoshal-Zamolodchikov {[}23{]}), which lies
outside the scope of this paper.

\begin{center}\rule{0.5\linewidth}{0.5pt}\end{center}

\subsection{10. Conclusion}\label{conclusion}

In this paper, we have provided an overview of the mathematical
structure of the sine-Gordon equation and demonstrated the following:

\begin{enumerate}
\def\labelenumi{\arabic{enumi}.}
\item
  Soliton solutions (kinks, anti-kinks, breathers) are
  \textbf{rigorously} characterized by a finite number of parameters
  \((N, \theta_0, \alpha)\) (§8).
\item
  All structural properties of solitons -- conservation of winding
  number, shape preservation during collisions and phase shift, finite
  parameter representation -- are preserved in the rectangular wave
  limit \(\sigma \to 0\) (§9).
\item
  The finiteness of energy is based on the structure in which the energy
  of each \(2\pi\) rotation is self-contained (the potential
  \(V(\phi) = 1-\cos\phi\) returns from \(0 \to 0\) over one period). A
  kink-anti-kink pair (\(4\pi\) rotation) can undergo pair annihilation,
  and the permanent energy is only the contribution from the remaining
  kinks that cannot undergo pair annihilation.
\item
  Therefore, the transition from solitons to rectangular waves involves
  no loss of information. Among the infinite degrees of freedom of the
  continuous field, the physically meaningful information is
  concentrated in the finite number of parameters
  \((N, \theta_0, \alpha)\), and this concentration holds regardless of
  the soliton shape.
\end{enumerate}

\textbf{Limits of applicability:} The above conclusions hold under the
implicit assumptions of this theory (Proposition 9.6a).

\begin{itemize}
\tightlist
\item
  \textbf{Spatial constraints}: The phase at the origin is fixed at the
  vacuum value \(\phi = 0\), and the phase at the terminus is forced to
  the vacuum value \(\phi = 2n\pi\) (the winding number \(N\) is
  integers only). Boundary reflections do not exist. These conditions
  are satisfied only on infinite space \((-\infty, +\infty)\) or on a
  closed curve \(S^1\). Extension to open strings (boundary conditions
  on a finite interval) requires admitting half-integer winding numbers
  (\(\phi = (2n+1)\pi\), potential maximum) and handling soliton
  reflection at boundaries, which lies outside the scope of this paper.
\item
  \textbf{Collision constraints}: This theory is a completely integrable
  system, and soliton collisions are described as perfectly elastic
  collisions. To avoid violating momentum conservation, the velocity,
  shape, and energy of each soliton are individually conserved before
  and after the collision, and only the position (phase shift
  \(\Delta x\)) changes. This effectively means that solitons behave as
  free particles that do not interact with each other, and dissipation
  of energy, exchange of velocities, and inelastic channels (pair
  annihilation → radiation) do not occur within the scope of classical
  theory.
\end{itemize}

The rectangular wave description is not a ``coarse approximation'' of
the soliton but rather a rigorous limit that completely inherits its
topological and scattering structure. The fact that solutions of
nonlinear wave equations are reduced to a finite number of discrete
parameters reveals a deep correspondence between continuous field
theories and discrete structures.

The above limits of applicability -- in particular, the inability to
handle systems with different inertia parameters and the absence of
velocity redistribution after collisions -- are overcome by the model
presented in §12. The definition of transformation invariants for the
sinusoidal-rectangular wave transformation is given in §11.

\begin{center}\rule{0.5\linewidth}{0.5pt}\end{center}

\subsection{11. Sinusoidal-Rectangular Wave Transformation and
Transformation
Invariants}\label{sinusoidal-rectangular-wave-transformation-and-transformation-invariants}

\subsubsection{11.1 Purpose}\label{purpose-1}

In §9, we showed that soliton solutions preserve their structure even in
the rectangular wave limit. In this section, we rigorously define the
\textbf{forward transform} and \textbf{inverse transform} between
sinusoidal and rectangular waveforms, and derive quantities invariant
under the transformation (transformation invariants).

The invariants defined in this section are dimensionless quantities, and
no identification with specific physical quantities is made.

\subsubsection{11.2 Definition of
Waveforms}\label{definition-of-waveforms}

We define periodic waveforms \(\phi(\theta)\) on phase space as
functions on the phase variable \(\theta \in [0, 2\pi)\). Here
\(\theta\) is an index of processing steps, not time or distance.

\textbf{Sinusoidal waveform:}

\[\phi_{\sin}(\theta) = A_{\sin} \cdot \sin(\theta) \tag{W1}\]

\textbf{Rectangular waveform:}

\[\phi_{\text{rect}}(\theta) = A_{\text{rect}} \cdot \mathrm{sgn}(\sin(\theta)) \tag{W2}\]

Here \(A_{\sin}\) and \(A_{\text{rect}}\) are the maximum wave heights
(amplitudes) of the respective waveforms.

\subsubsection{11.3 Definition of Waveform
Coefficient}\label{definition-of-waveform-coefficient}

\textbf{Definition 11.1 (Half-period area).} The half-period area
\(S[\phi]\) of a waveform \(\phi(\theta)\) is defined as the integral of
the absolute value of the waveform over a half-period (phase width
\(\pi\)):

\[S[\phi] = \int_0^{\pi} |\phi(\theta)|\, d\theta \tag{W3}\]

We compute the half-period area for each waveform.

\textbf{Sinusoidal wave:}

\[S[\phi_{\sin}] = \int_0^{\pi} A_{\sin} \sin(\theta)\, d\theta = A_{\sin} \left[-\cos(\theta)\right]_0^{\pi} = 2A_{\sin} \tag{W4}\]

\textbf{Rectangular wave:}

\[S[\phi_{\text{rect}}] = \int_0^{\pi} A_{\text{rect}} \cdot 1\, d\theta = \pi \cdot A_{\text{rect}} \tag{W5}\]

\textbf{Definition 11.2 (Waveform coefficient).} The amplitude
transformation coefficient \(\kappa\) is defined as the factor that
equalizes the half-period areas of the sinusoidal and rectangular waves:

\[\kappa = \frac{2}{\pi} \approx 0.6366 \tag{W6}\]

\subsubsection{11.4 Forward Transform and Inverse
Transform}\label{forward-transform-and-inverse-transform}

\textbf{Definition 11.3 (Forward transform \(T\)).} The transform \(T\)
from a sinusoidal waveform to a rectangular waveform is defined as
follows:

\[T: \phi_{\sin}(\theta) \mapsto \phi_{\text{rect}}(\theta), \qquad A_{\text{rect}} = \kappa \cdot A_{\sin} = \frac{2}{\pi} A_{\sin} \tag{W7}\]

That is, the waveform is rectangularized while simultaneously reducing
the maximum wave height by a factor of \(\kappa\).

\textbf{Definition 11.4 (Inverse transform \(T^{-1}\)).} The transform
\(T^{-1}\) from a rectangular waveform to a sinusoidal waveform is
defined as follows:

\[T^{-1}: \phi_{\text{rect}}(\theta) \mapsto \phi_{\sin}(\theta), \qquad A_{\sin} = \frac{1}{\kappa} \cdot A_{\text{rect}} = \frac{\pi}{2} A_{\text{rect}} \tag{W8}\]

That is, the waveform is sinusoidalized while simultaneously increasing
the maximum wave height by a factor of \(1/\kappa\).

\textbf{Proposition 11.1 (Existence of the inverse).}
\(T^{-1} \circ T = T \circ T^{-1} = \mathrm{id}\) (identity
transformation) holds.

\emph{Proof.} The amplitude of \(T^{-1}(T(\phi_{\sin}))\) is
\((\pi/2) \cdot (2/\pi) \cdot A_{\sin} = A_{\sin}\), and the waveform
returns to a sinusoidal wave. \(\square\)

\subsubsection{\texorpdfstring{11.5 Transformation Invariant \(I_1\)
(Quantitative
Invariant)}{11.5 Transformation Invariant I\_1 (Quantitative Invariant)}}\label{transformation-invariant-i_1-quantitative-invariant}

\textbf{Definition 11.5.} The quantitative invariant \(I_1[\phi]\) of a
waveform \(\phi\) is defined as the half-period area:

\[I_1[\phi] = S[\phi] = \int_0^{\pi} |\phi(\theta)|\, d\theta \tag{W9}\]

\textbf{Theorem 11.1 (Transformation invariance of \(I_1\)).} The
quantitative invariant \(I_1\) is invariant under the forward transform
\(T\) and the inverse transform \(T^{-1}\):

\[I_1[T(\phi_{\sin})] = I_1[\phi_{\sin}]\]

\emph{Proof.}

\[I_1[\phi_{\sin}] = 2A_{\sin} \qquad \text{(from W4)}\]

\[I_1[T(\phi_{\sin})] = \pi \cdot A_{\text{rect}} = \pi \cdot \frac{2}{\pi} A_{\sin} = 2A_{\sin} \qquad \text{(from W5, W7)}\]

Both are equal. \(\square\)

\subsubsection{\texorpdfstring{11.6 Transformation Invariant \(I_2\)
(Translational
Invariant)}{11.6 Transformation Invariant I\_2 (Translational Invariant)}}\label{transformation-invariant-i_2-translational-invariant}

\(I_1\) is the half-period area of the waveform (Definition 11.1). As an
invariant expressing how much this quantity ``moves'' when the phase is
displaced by one step, we define \(I_2\). Through the introduction of a
normalization coefficient \(\kappa_\Delta\), the value of \(I_2\) is
preserved under waveform transformation.

\textbf{Definition 11.6 (Phase displacement per processing step).} The
phase displacement \(\Delta\phi\) per processing step (phase increment
\(\delta\)) of a waveform \(\phi(\theta)\) is defined as:

\[\Delta\phi = \phi(\theta + \delta) - \phi(\theta) \tag{W10}\]

\textbf{Definition 11.7 (Translational invariant).} The translational
invariant \(I_2[\phi]\) of a waveform \(\phi\) is defined as the product
of the phase displacement and the quantitative invariant:

\[I_2[\phi] = \kappa_\Delta \cdot \Delta\phi \cdot I_1[\phi] \tag{W11}\]

Here \(\kappa_\Delta\) is a normalization coefficient that depends on
the waveform, defined using the same waveform coefficient
\(\kappa = 2/\pi\) as for \(I_1\):

\begin{itemize}
\tightlist
\item
  Sinusoidal wave: \(\kappa_\Delta = 1\)
\item
  Rectangular wave: \(\kappa_\Delta = 1/\kappa = \pi/2\)
\end{itemize}

\textbf{Theorem 11.2 (Transformation invariance of \(I_2\)).} The
translational invariant \(I_2\) is invariant under the forward transform
\(T\) and the inverse transform \(T^{-1}\).

\emph{Proof.} The forward transform \(T\) multiplies the amplitude by
\(\kappa\) and multiplies the normalization coefficient by \(1/\kappa\).
\(I_1\) is invariant by Theorem 11.1. Since the phase displacement is
proportional to the amplitude, it is multiplied by \(\kappa\), but
\(\kappa_\Delta\) is multiplied by \(1/\kappa\), so the product
\(\kappa_\Delta \cdot \Delta\phi\) is invariant.

\[I_2[T(\phi)] = \frac{1}{\kappa} \cdot (\kappa \cdot \Delta\phi) \cdot I_1 = \Delta\phi \cdot I_1 = I_2[\phi] \qquad \square\]

\subsubsection{11.7 Summary of Transformation
Invariants}\label{summary-of-transformation-invariants}

\textbf{Theorem 11.3 (Invariants of the sinusoidal-rectangular wave
transformation).} Under the forward transform \(T\) (Definition 11.3)
and inverse transform \(T^{-1}\) (Definition 11.4) using the waveform
coefficient \(\kappa = 2/\pi\), the following quantities are invariant:

{\def\LTcaptype{none} % do not increment counter
\begin{longtable}[]{@{}
  >{\centering\arraybackslash}p{(\linewidth - 6\tabcolsep) * \real{0.3333}}
  >{\raggedright\arraybackslash}p{(\linewidth - 6\tabcolsep) * \real{0.2500}}
  >{\raggedright\arraybackslash}p{(\linewidth - 6\tabcolsep) * \real{0.1667}}
  >{\raggedright\arraybackslash}p{(\linewidth - 6\tabcolsep) * \real{0.2500}}@{}}
\toprule\noalign{}
\begin{minipage}[b]{\linewidth}\centering
Invariant
\end{minipage} & \begin{minipage}[b]{\linewidth}\raggedright
Definition
\end{minipage} & \begin{minipage}[b]{\linewidth}\raggedright
Value
\end{minipage} & \begin{minipage}[b]{\linewidth}\raggedright
Property
\end{minipage} \\
\midrule\noalign{}
\endhead
\bottomrule\noalign{}
\endlastfoot
\(I_0\) & Winding number \(N = [\phi(2\pi) - \phi(0)] / 2\pi\) & Integer
& Topological invariant \\
\(I_1\) & Half-period area \(\int_0^{\pi} |\phi(\theta)|\, d\theta\) &
\(2A_{\sin}\) & Positive definite, additive \\
\(I_2\) & \(\kappa_\Delta \cdot \Delta\phi \cdot I_1\) &
\(\Delta\phi \cdot 2A_{\sin}\) & Signed, additive \\
\end{longtable}
}

\emph{Proof.} \(I_0\): depends only on boundary values and is
independent of waveform (Proposition 9.1). \(I_1\): Theorem 11.1.
\(I_2\): Theorem 11.2. All are transformation invariant due to the
cancellation of the waveform coefficient \(\kappa\). \(\square\)

\textbf{Corollary 11.1.} The rectangular wave representation is not a
``coarse approximation'' of the sinusoidal wave representation. The
transform \(T\) is a bijection that preserves the three transformation
invariants \(I_0, I_1, I_2\). The values of the invariants computed on
the rectangular wave coincide exactly with the values on the
corresponding sinusoidal wave.

\textbf{Corollary 11.2.} The invariants are determined solely from the
waveform parameters (amplitude \(A\), winding number \(N\)) and do not
depend on the internal structure of the waveform (whether sinusoidal or
rectangular).

\begin{center}\rule{0.5\linewidth}{0.5pt}\end{center}

\subsection{12. Elastic State Transformation Model on Phase
Space}\label{elastic-state-transformation-model-on-phase-space}

In §9--10, we identified the limits of applicability of the sine-Gordon
soliton theory. In particular, (i) the requirement that all solitons
have the same inertia parameter, and (ii) the absence of velocity
redistribution after collisions (transmissive collisions), are
insufficient for describing general elastic state transformations.

In this section, as an alternative model without these constraints, we
formulate \textbf{elastic state transformations at arbitrary phases for
oscillating systems with different inertia parameters}, predicated on
the preservation of the transformation invariants \(I_1, I_2\) defined
in §11. All parameters of this model are abstract quantities on phase
space, and no identification with specific physical quantities is made.

\subsubsection{12.1 Definition of the Oscillating
System}\label{definition-of-the-oscillating-system}

We define a harmonic oscillating system on phase space. The equation of
motion for a system with inertia parameter \(M\) and restoring
coefficient \(k\) is

\[M\ddot{x} + 2kx = 0 \tag{W12}\]

Here \(x\) is the displacement from the equilibrium position, and
\(\dot{x}\) is the rate of change of displacement per processing step.
The system has angular frequency \(\omega = \sqrt{2k/M}\), and the
general solution is

\[x(\theta) = A\cos(\omega\theta + \phi_0) \tag{W13}\]

Here \(A\) is the amplitude, \(\phi_0\) is the initial phase, and
\(\theta\) is the index of processing steps.

\subsubsection{12.2 Configuration of the Two-Body System and Definition
of State
Transformation}\label{configuration-of-the-two-body-system-and-definition-of-state-transformation}

Consider two independent oscillating systems with different inertia
parameters \(M_1, M_2\):

\[\ddot{x}_1 + \omega_1^2 x_1 = 0, \qquad \omega_1 = \sqrt{\frac{2k}{M_1}} \tag{W14}\]

\[\ddot{x}_2 + \omega_2^2 x_2 = 0, \qquad \omega_2 = \sqrt{\frac{2k}{M_2}} \tag{W15}\]

When \(M_1 \neq M_2\), we have \(\omega_1 \neq \omega_2\), and the two
systems oscillate independently.

\textbf{Definition of state transformation.} In this paper, a ``state
transformation'' is defined as a mathematical process in which the
states \((x_1, \dot{x}_1)\) and \((x_2, \dot{x}_2)\) of both systems are
coupled through two conservation conditions at a certain step
\(\theta = \theta_c\), instantaneously transitioning to new states
\((x_1, \dot{x}_1')\) and \((x_2, \dot{x}_2')\). The positions \(x_i\)
do not change through the state transformation; only the rates of change
\(v_i = \dot{x}_i\) change according to (W20)--(W21).

\textbf{Remark.} The state transformation is defined as an abstract
operation: ``an instantaneous state transition in which the states of
two oscillating systems are coupled through conservation conditions.''
The step \(\theta_c\) at which the state transformation occurs is an
input parameter given externally, and is not determined from within the
model. This allows state transformations at arbitrary phases
\(\phi_1, \phi_2\) to be treated in a unified manner.

\subsubsection{12.3 Elastic State Transformation at Arbitrary
Phases}\label{elastic-state-transformation-at-arbitrary-phases}

\textbf{State immediately before the state transformation:} The state of
each system at step \(\theta_c\) is expressed as:

\[x_{1,c} = A_1\cos\phi_1, \qquad v_1 = -A_1\omega_1\sin\phi_1 \tag{W16}\]

\[x_{2,c} = A_2\cos\phi_2, \qquad v_2 = -A_2\omega_2\sin\phi_2 \tag{W17}\]

Here \(\phi_i = \omega_i\theta_c + \phi_{0,i}\) is the phase of each
system at the state transformation step, and \(A_i\) is the amplitude.

\textbf{Conservation conditions for elastic state transformation:}

\[M_1 v_1 + M_2 v_2 = M_1 v_1' + M_2 v_2' \tag{W18}\]

\[\frac{1}{2}M_1 v_1^2 + \frac{1}{2}M_2 v_2^2 = \frac{1}{2}M_1 v_1'^2 + \frac{1}{2}M_2 v_2'^2 \tag{W19}\]

Equation (W18) represents the conservation of the first-order form of
the rate of change \(\sum M_i v_i\), which has the same structure as
\(I_2\) (translational invariant: area \(\times\) displacement) of
§11.6. Equation (W19) represents the conservation of the second-order
form of the rate of change \(\sum \frac{1}{2}M_i v_i^2\), which has the
same structure as \(I_1\) (quantitative invariant: positive-definite
additive conserved quantity) of §11.5. Both are mathematically
isomorphic conservation conditions, and from these two conditions, the
rates of change after the state transformation are uniquely determined.

\textbf{Theorem 12.1 (Elastic state transformation formula at arbitrary
phases).} The rates of change after transformation for systems with
inertia parameters \(M_1, M_2\) undergoing state transformation at
phases \(\phi_1, \phi_2\) are given by:

\[v_1' = \frac{(M_1 - M_2)\,v_1 + 2M_2\,v_2}{M_1 + M_2} \tag{W20}\]

\[v_2' = \frac{(M_2 - M_1)\,v_2 + 2M_1\,v_1}{M_1 + M_2} \tag{W21}\]

\emph{Proof.} Obtained immediately by solving the simultaneous equations
(W18) and (W19). \(\square\)

\subsubsection{12.4 Oscillation State After State
Transformation}\label{oscillation-state-after-state-transformation}

The state transformation is instantaneous, and positions do not change.
After the state transformation, each system resumes oscillation with the
position \(x_{i,c}\) at the time of transformation and the new rate of
change \(v_i'\) as initial conditions:

\[x_1(\theta) = x_{1,c}\cos(\omega_1\tau) + \frac{v_1'}{\omega_1}\sin(\omega_1\tau) \qquad (\tau = \theta - \theta_c) \tag{W22}\]

\[x_2(\theta) = x_{2,c}\cos(\omega_2\tau) + \frac{v_2'}{\omega_2}\sin(\omega_2\tau) \tag{W23}\]

\textbf{Amplitude and phase after state transformation:}

\[A_i' = \sqrt{x_{i,c}^2 + \left(\frac{v_i'}{\omega_i}\right)^2} \tag{W24}\]

\[\phi_i' = \arctan\!\left(\frac{-v_i'}{\omega_i\, x_{i,c}}\right) \tag{W25}\]

The angular frequency \(\omega_i\) is determined solely by the restoring
coefficient and inertia parameter, and therefore does not change through
the state transformation. What changes are only the amplitude \(A_i\)
and the phase \(\phi_i\).

\subsubsection{\texorpdfstring{12.5 Proof of \(I_1\)
Conservation}{12.5 Proof of I\_1 Conservation}}\label{proof-of-i_1-conservation}

\textbf{Proposition 12.1.} In an elastic state transformation at
arbitrary phases, the total quantitative invariant of the system is
conserved.

\emph{Proof.} The total conserved quantity for each system is the sum of
a position-dependent term and a rate-of-change-dependent term:

\[E_i = \frac{1}{2}M_i\omega_i^2\,x_{i,c}^2 + \frac{1}{2}M_i\,v_i^2 = k\,x_{i,c}^2 + \frac{1}{2}M_i\,v_i^2\]

Since the state transformation is instantaneous, the position-dependent
term \(k\,x_{i,c}^2\) is invariant before and after the transformation.
The rate-of-change-dependent term is conserved by the conservation
condition (W19). Therefore

\[E_1 + E_2 = E_1' + E_2' \qquad \square \tag{W26}\]

\subsubsection{12.6 Characteristics of Amplitude and Phase Changes by
State Transformation
Phase}\label{characteristics-of-amplitude-and-phase-changes-by-state-transformation-phase}

From the state transformation formulas (W20)--(W21) and the
post-transformation amplitude and phase determination formulas
(W24)(W25), the following characteristics appear depending on the phase
\(\phi_i\) at the transformation step:

{\def\LTcaptype{none} % do not increment counter
\begin{longtable}[]{@{}
  >{\raggedright\arraybackslash}p{(\linewidth - 8\tabcolsep) * \real{0.2000}}
  >{\raggedright\arraybackslash}p{(\linewidth - 8\tabcolsep) * \real{0.2000}}
  >{\raggedright\arraybackslash}p{(\linewidth - 8\tabcolsep) * \real{0.2000}}
  >{\raggedright\arraybackslash}p{(\linewidth - 8\tabcolsep) * \real{0.2000}}
  >{\raggedright\arraybackslash}p{(\linewidth - 8\tabcolsep) * \real{0.2000}}@{}}
\toprule\noalign{}
\begin{minipage}[b]{\linewidth}\raggedright
Transformation Phase
\end{minipage} & \begin{minipage}[b]{\linewidth}\raggedright
Position \(x_{i,c}\)
\end{minipage} & \begin{minipage}[b]{\linewidth}\raggedright
Rate of Change \(|v_i|\)
\end{minipage} & \begin{minipage}[b]{\linewidth}\raggedright
Conserved Quantity Distribution
\end{minipage} & \begin{minipage}[b]{\linewidth}\raggedright
Effect of State Transformation
\end{minipage} \\
\midrule\noalign{}
\endhead
\bottomrule\noalign{}
\endlastfoot
\(\phi_i = \pi/2\) & \(0\) & \(A_i\omega_i\) (maximum) & Entirely in
rate-of-change-dependent terms & Rates of change are fully
redistributed; only amplitudes \(A_i\) change \\
\(\phi_i = 0\) & \(A_i\) (maximum) & \(0\) & Entirely in
position-dependent terms & No exchange of rates of change occurs; state
is invariant before and after \\
\(\phi_i = \pi/4\) & \(A_i/\sqrt{2}\) & \(A_i\omega_i/\sqrt{2}\) & Equal
distribution between position and rate of change & Partial
redistribution of rates of change and phase shift occur \\
\end{longtable}
}

\textbf{Remark.} The ``invariance of state'' at \(\phi_i = 0\) is
because the right-hand sides of (W20)(W21) are determined solely by
\(v_i\), and when \(v_i = 0\), \(v_i' = 0\) is the solution. In the
general case \(\phi_1 \neq \phi_2\), as long as the rates of change of
both systems are not simultaneously zero, a non-trivial state
transformation necessarily occurs.

\subsubsection{12.7 Comparison with sine-Gordon
Theory}\label{comparison-with-sine-gordon-theory}

{\def\LTcaptype{none} % do not increment counter
\begin{longtable}[]{@{}
  >{\raggedright\arraybackslash}p{(\linewidth - 4\tabcolsep) * \real{0.3333}}
  >{\raggedright\arraybackslash}p{(\linewidth - 4\tabcolsep) * \real{0.3333}}
  >{\raggedright\arraybackslash}p{(\linewidth - 4\tabcolsep) * \real{0.3333}}@{}}
\toprule\noalign{}
\begin{minipage}[b]{\linewidth}\raggedright
Property
\end{minipage} & \begin{minipage}[b]{\linewidth}\raggedright
sine-Gordon (§1--10)
\end{minipage} & \begin{minipage}[b]{\linewidth}\raggedright
This Model (§12)
\end{minipage} \\
\midrule\noalign{}
\endhead
\bottomrule\noalign{}
\endlastfoot
Different inertia parameters & Not possible (all kinks identical) &
\textbf{Possible} (\(M_1 \neq M_2\)) \\
Rate-of-change redistribution after transformation & None (transmissive)
& \textbf{Present} (state transformation formula) \\
State transformation at arbitrary phases & Undefined & \textbf{Fully
defined} \\
Closed string assumption & Required (implicit) & \textbf{Not
required} \\
Integer constraint on winding number & Required & \textbf{Not
required} \\
Half-integer phases & Excluded (\(\phi = \pi\) not allowed) &
\textbf{Allowed} \\
Wave equation & \(\phi_{tt} - \phi_{xx} + \sin\phi = 0\) &
\(\ddot{x}_i + \omega_i^2 x_i = 0\) \\
Hidden assumptions & 6 (Proposition 9.6a) & \textbf{None} (conservation
conditions only) \\
Dissipation & Zero & Zero \\
Premises of formulation & Conservation of the first-order form
(\(\simeq I_2\)) + conservation of the second-order form
(\(\simeq I_1\)) + topological constraint + integrability + closed space
+ identical inertia & \textbf{\(I_2\) conservation + \(I_1\)
conservation only} \\
\end{longtable}
}

\subsubsection{12.8 Equivalence with the Rectangular Wave
Limit}\label{equivalence-with-the-rectangular-wave-limit}

In §9, we showed that solitons of the sine-Gordon equation preserve
their structure in the rectangular wave limit. In this section, we prove
that in the model of §12 as well, the results of state transformation
are identical regardless of whether the waveform changes from sinusoidal
(cosine) to rectangular.

\textbf{Introduction of waveform parameters.} We define the oscillation
waveform as a family \(x^\sigma(\theta)\) controlled by a transition
width parameter \(\sigma\) (\(0 < \sigma \leq 1\)):

\begin{itemize}
\tightlist
\item
  \(\sigma = 1\): smooth sinusoidal (cosine) wave
  \(x^1(\theta) = A\cos(\omega\theta)\)
\item
  \(\sigma \to 0\): rectangular wave
  \(x^0(\theta) = A\,\mathrm{sgn}(\cos(\omega\theta))\)
\end{itemize}

Intermediate \(\sigma\) controls the steepness of the transition region.
For all \(\sigma\), the amplitude is \(A\) and the period is
\(T = 2\pi/\omega\) in common.

\textbf{Theorem 12.2 (Waveform independence of state transformation
results).} When the position \(x_{i,c}\) and total conserved quantity
\(E_i\) at the time of state transformation are given, the
post-transformation rate of change \(v_i'\), amplitude \(A_i'\), and
phase \(\phi_i'\) do not depend on the waveform parameter \(\sigma\).

\emph{Proof.} We show this in three stages.

\textbf{Stage 1: The rate of change at the time of transformation is
uniquely determined from the conserved quantity and position.}

The total conserved quantity of the system is preserved regardless of
the waveform:

\[E_i = k\,x_i^2(\theta) + \frac{1}{2}M_i\,\dot{x}_i^2(\theta) = k\,A_i^2 \tag{W27}\]

This does not depend on the waveform. The position-dependent term is a
function of position only, and the rate-of-change-dependent term is the
remainder. Therefore, the rate of change at the transformation position
\(x_{i,c}\) is

\[v_i^2 = \frac{2(E_i - k\,x_{i,c}^2)}{M_i} = \omega_i^2(A_i^2 - x_{i,c}^2) \tag{W28}\]

and is \textbf{uniquely} determined. This does not depend on \(\sigma\).
\(\checkmark\)

\textbf{Stage 2: The state transformation formula depends only on the
rates of change.}

The state transformation formulas (W20)--(W21) of Theorem 12.1 are
algebraic relations involving only \(M_1, M_2, v_1, v_2\), and contain
no waveform information. Since Stage 1 showed that \(v_i\) is
\(\sigma\)-independent, \(v_i'\) is also \(\sigma\)-independent.
\(\checkmark\)

\textbf{Stage 3: The post-transformation amplitude and phase are also
\(\sigma\)-independent.}

The total conserved quantity after transformation is

\[E_i' = k\,x_{i,c}^2 + \frac{1}{2}M_i\,v_i'^2 \tag{W29}\]

and the new amplitude is

\[A_i' = \sqrt{\frac{E_i'}{k}} = \sqrt{x_{i,c}^2 + \frac{v_i'^2}{\omega_i^2}} \tag{W30}\]

Since all of \(x_{i,c}\), \(v_i'\), \(\omega_i\) are
\(\sigma\)-independent, \(A_i'\) is also \(\sigma\)-independent.

The initial phase \(\phi_i' = \arctan(-v_i'/(\omega_i\,x_{i,c}))\) is
likewise \(\sigma\)-independent. \(\square\)

\textbf{Corollary 12.1 (Rigorous nature of the rectangular wave limit).}
The results of elastic state transformation in this model are
\textbf{exactly identical} whether the oscillation is a smooth
sinusoidal wave or a rectangular wave. The only information that is lost
is the detail of the waveform (the shape of the transition region),
while the redistribution of rates of change, conservation of conserved
quantities, and changes in amplitude and phase do not depend on the
waveform at all.

\textbf{Corollary 12.2.} This result reproduces, in a more general
framework, the structural preservation of the rectangular wave limit
demonstrated for the sine-Gordon equation in §9. While the proof of
structural preservation in the sine-Gordon case required the assumptions
of integrability and closed space, in this model the equivalence holds
with \textbf{only the single condition of \(I_1\) conservation}
(equation (W27)).

\subsubsection{12.9 Verification Summary}\label{verification-summary}

{\def\LTcaptype{none} % do not increment counter
\begin{longtable}[]{@{}
  >{\centering\arraybackslash}p{(\linewidth - 8\tabcolsep) * \real{0.1364}}
  >{\centering\arraybackslash}p{(\linewidth - 8\tabcolsep) * \real{0.1818}}
  >{\centering\arraybackslash}p{(\linewidth - 8\tabcolsep) * \real{0.1818}}
  >{\centering\arraybackslash}p{(\linewidth - 8\tabcolsep) * \real{0.3636}}
  >{\centering\arraybackslash}p{(\linewidth - 8\tabcolsep) * \real{0.1364}}@{}}
\toprule\noalign{}
\begin{minipage}[b]{\linewidth}\centering
Property
\end{minipage} & \begin{minipage}[b]{\linewidth}\centering
Sinusoidal
\end{minipage} & \begin{minipage}[b]{\linewidth}\centering
Rectangular
\end{minipage} & \begin{minipage}[b]{\linewidth}\centering
\(\sigma\)-dependence
\end{minipage} & \begin{minipage}[b]{\linewidth}\centering
Remarks
\end{minipage} \\
\midrule\noalign{}
\endhead
\bottomrule\noalign{}
\endlastfoot
Rate of change at transformation \(v_i\) & \(\checkmark\) &
\(\checkmark\) & \textbf{Independent} & By (W28) \\
Post-transformation rate of change \(v_i'\) & \(\checkmark\) &
\(\checkmark\) & \textbf{Independent} & (W20)--(W21) \\
Post-transformation amplitude \(A_i'\) & \(\checkmark\) & \(\checkmark\)
& \textbf{Independent} & (W30) \\
Post-transformation phase \(\phi_i'\) & \(\checkmark\) & \(\checkmark\)
& \textbf{Independent} & (W25) \\
\(I_1\) conservation & \(\checkmark\) & \(\checkmark\) &
\textbf{Independent} & (W26) \\
\(I_2\) conservation & \(\checkmark\) & \(\checkmark\) &
\textbf{Independent} & (W18) \\
Different inertia parameters & \(\checkmark\) & \(\checkmark\) &
\textbf{Independent} & \(M_1 \neq M_2\) \\
Arbitrary phases & \(\checkmark\) & \(\checkmark\) &
\textbf{Independent} & \(\phi_1, \phi_2\) arbitrary \\
\end{longtable}
}

\textbf{Theorem 12.3 (Structural preservation of the elastic state
transformation model).} The process in which two oscillating systems
with different inertia parameters \(M_1 \neq M_2\) undergo elastic state
transformation at arbitrary phases \(\phi_1, \phi_2\) is completely
described by the wave equations (W14)--(W15) and the state
transformation formulas (W20)--(W21), and all of the following
properties are preserved in the rectangular wave limit \(\sigma \to 0\):

\begin{enumerate}
\def\labelenumi{\arabic{enumi}.}
\tightlist
\item
  \textbf{Redistribution of rates of change} (state transformation
  formula: \(\sigma\)-independent)
\item
  \textbf{Conservation of \(I_1\) and \(I_2\)} ((W26) and (W18):
  \(\sigma\)-independent)
\item
  \textbf{Determination of amplitude and phase} ((W30) and (W25):
  \(\sigma\)-independent)
\end{enumerate}

\emph{Proof.} A direct consequence of Theorem 12.2. \(\square\)

The only assumptions of this formulation are the following two
conservation conditions:

\begin{itemize}
\tightlist
\item
  \textbf{\(I_2\) conservation}:
  \(M_1 v_1 + M_2 v_2 = M_1 v_1' + M_2 v_2'\)
\item
  \textbf{\(I_1\) conservation}:
  \(\frac{1}{2}M_1 v_1^2 + \frac{1}{2}M_2 v_2^2 = \frac{1}{2}M_1 v_1'^2 + \frac{1}{2}M_2 v_2'^2\)
\end{itemize}

Topological constraints, integrability, compactness of space,
integrality of winding numbers, periodicity of the potential, and the
assumption of identical inertia parameters are all unnecessary, and the
equivalence with the rectangular wave limit follows automatically from
\(I_1\) conservation.

\begin{center}\rule{0.5\linewidth}{0.5pt}\end{center}

\subsection{References}\label{references}

{[}1{]} Frenkel, J. \& Kontorova, T. ``On the theory of plastic
deformation and twinning.'' \emph{Izv. Akad. Nauk SSSR, Ser. Fiz.} 1,
137--149 (1939).

{[}2{]} Scott, A.C., Chu, F.Y.F. \& McLaughlin, D.W. ``The Soliton: A
New Concept in Applied Science.'' \emph{Proc. IEEE} 61, 1443--1483
(1973).

{[}3{]} Rubinstein, J. ``Sine-Gordon equation.'' \emph{J. Math. Phys.}
11, 258--266 (1970).

{[}4{]} Rajaraman, R. \emph{Solitons and Instantons: An Introduction to
Solitons and Instantons in Quantum Field Theory.} North-Holland (1982).

{[}5{]} Bour, E. ``Théorie de la déformation des surfaces.'' \emph{J.
Éc. Imp. Polytech.} 22, Cahier 39, 1--148 (1862).

{[}6{]} Enneper, A. ``Über asymptotische Linien.'' \emph{Nachr. Königl.
Gesellsch. d.~Wiss. Göttingen} 1870, 493--510 (1870).

{[}7{]} Bianchi, L. ``Ricerche sulle superficie a curvatura costante e
sulle elicoidi.'' \emph{Ann. Sc. Norm. Super. Pisa} 2, 285--341 (1879).

{[}8{]} Skyrme, T.H.R. ``A non-linear field theory.'' \emph{Proc. R.
Soc. Lond. A} 260, 127--138 (1961).

{[}9{]} Perring, J.K. \& Skyrme, T.H.R. ``A model unified field
equation.'' \emph{Nucl. Phys.} 31, 550--555 (1962).

{[}10{]} Manton, N.S. \& Sutcliffe, P.M. \emph{Topological Solitons.}
Cambridge University Press (2004).

{[}11{]} Ablowitz, M.J., Kaup, D.J., Newell, A.C. \& Segur, H. ``Method
for solving the sine-Gordon equation.'' \emph{Phys. Rev.~Lett.} 30,
1262--1264 (1973).

{[}12{]} Ablowitz, M.J., Kaup, D.J., Newell, A.C. \& Segur, H. ``The
inverse scattering transform --- Fourier analysis for nonlinear
problems.'' \emph{Stud. Appl. Math.} 53, 249--315 (1974).

{[}13{]} Faddeev, L.D. \& Takhtajan, L.A. \emph{Hamiltonian Methods in
the Theory of Solitons.} Springer (1987).

{[}14{]} Zakharov, V.E. \& Shabat, A.B. ``Exact theory of
two-dimensional self-focusing and one-dimensional self-modulation of
waves in nonlinear media.'' \emph{Sov. Phys. JETP} 34, 62--69 (1972).

{[}15{]} Rogers, C. \& Schief, W.K. \emph{Bäcklund and Darboux
Transformations: Geometry and Modern Applications in Soliton Theory.}
Cambridge University Press (2002).

{[}16{]} Gibbon, J.D., James, I.N. \& Moroz, I.M. ``The sine-Gordon
equation as a model for a rapidly rotating baroclinic fluid.''
\emph{Proc. R. Soc. Lond. A} 367, 219--237 (1979).

{[}17{]} Derrick, G.H. ``Comments on nonlinear wave equations as models
for elementary particles.'' \emph{J. Math. Phys.} 5, 1252--1254 (1964).

{[}18{]} Manton, N.S. ``Geometry of Skyrmions.'' \emph{Commun. Math.
Phys.} 111, 469--478 (1987).

{[}19{]} Barone, A. \& Paternò, G. \emph{Physics and Applications of the
Josephson Effect.} Wiley (1982).

{[}20{]} Coleman, S. ``Quantum sine-Gordon equation as the massive
Thirring model.'' \emph{Phys. Rev.~D} 11, 2088--2097 (1975).

{[}21{]} Mandelstam, S. ``Soliton operators for the quantized
sine-Gordon equation.'' \emph{Phys. Rev.~D} 11, 3026--3030 (1975).

{[}22{]} Manton, N.S. ``A remark on the scattering of BPS monopoles.''
\emph{Phys. Lett. B} 110, 54--56 (1982).

{[}23{]} Ghoshal, S. \& Zamolodchikov, A.B. ``Boundary S matrix and
boundary state in two-dimensional integrable quantum field theory.''
\emph{Int. J. Mod. Phys. A} 9, 3841--3885 (1994).

\end{document}
